\documentclass[12pt]{article}
\usepackage[utf8]{inputenc}
\usepackage[T1]{fontenc}
\usepackage[margin=1in]{geometry}
\usepackage{setspace}
\usepackage{xcolor}
\usepackage[colorlinks=true, urlcolor=blue!60!black, linkcolor=blue!60!black]{hyperref}
\setlength{\parindent}{0em}
\setlength{\parskip}{0.45em}
\renewcommand{\baselinestretch}{1.0}
\begin{document}
\begin{center}
{\Large \textit{U.S. Department of Agriculture v. Moreno}, 413 U.S. 528 (1973)\par}
\vspace{0.3em}
{\normalsize Case Brief --- One-Page Edition\par}
\end{center}

\textbf{Memory Jogger.} Animus toward ``hippie communes'' cannot sustain a food stamp exclusion (\href{https://supreme.justia.com/cases/federal/us/413/528/}{Justia}).

\textbf{Facts.} A 1971 amendment barred food stamps for any household with an unrelated member, aiming at ``hippies.'' It cut off destitute households like that of Jacinta Moreno, a diabetic who lived with an unrelated family to survive.

\textbf{Procedural History.} A three-judge district court held the exclusion violated equal protection.

\textbf{Votes.} Affirmed 7--2: Brennan for the Court; Douglas concurred; Rehnquist, with Burger, dissented.

\textbf{Holding.} The classification fails rational-basis review: its purpose is illegitimate and its fraud rationale baseless.

\textbf{Reasoning.} Family ties are irrelevant to nutrition. The Act already punished fraud, abusers could evade the rule by splitting households, and only the most desperate were excluded. Douglas relied on freedom of association; Rehnquist said Congress could rationally fund only related households. \textit{Cleburne}, \textit{Romer}, and \textit{Windsor} later built on the anti-animus rule.

\textbf{Examples.} \emph{Same-Side:} a voucher ban on protesters. \emph{Opposite-Side:} a neutral zoning cap (\textit{Belle Terre}). \emph{Fence-Sitter:} a vehicle-dwelling ban.

\textbf{Critique.} \textit{Moreno} uses heightened scrutiny under a rational-basis label, and proving legislative animus from stray floor remarks raises hard questions of motive.

\textbf{Key Quotation.} ``For if the constitutional conception of `equal protection of the laws' means anything, it must, at the very least, mean that a bare congressional desire to harm a politically unpopular group cannot constitute a legitimate governmental interest.'' --- Justice Brennan, 413 U.S. at 534.

\end{document}
